\begin{question}
\noindent A historian is studying an old triangular star map recovered from an archive documenting early space observations. The triangular plot connects three observation points, $P$, $Q$ and $R$, marked on the map.

\noindent $PQ = 14.6$ cm, $QR = 10.3$ cm, and the angle at $R$ (angle $PRQ$) is $97^{\circ}$.

\begin{center}
\begin{tikzpicture}[scale=1, transform shape]
    \coordinate (R) at (-2,-1);
    \coordinate (Q) at (3,-1.3);
    \coordinate (P) at (1.3,4);

    \draw (P) -- (Q) -- (R) -- cycle;

    \pic [draw, -, angle radius=0.5cm, "$97^{\circ}$", angle eccentricity=1.6] {angle=Q--R--P};
    \pic [draw, -, angle radius=0.5cm, "$q^{\circ}$", angle eccentricity=1.6] {angle=R--P--Q};

    \node at (P) [above,font=\small] {$P$};
    \node at (Q) [right,font=\small] {$Q$};
    \node at (R) [left,font=\small] {$R$};

    \draw (P) -- (Q) node[midway, right] {$14.6$ cm};
    \draw (Q) -- (R) node[midway, below] {$10.3$ cm};
    \draw (R) -- (P) node[midway, above left] {};
\end{tikzpicture}
\end{center}

\noindent Work out the size of angle $QPR$, marked $q^{\circ}$ on the diagram.
Give your answer correct to 1 decimal place.

\vspace{5cm}

\begin{flushright}
    \answerequnit{q}{$^\circ$}{4}
\end{flushright}
\end{question}